% 拉普拉斯算子（综述）
% license CCBYSA3
% type Wiki

本文根据 CC-BY-SA 协议转载翻译自维基百科\href{https://en.wikipedia.org/wiki/Laplace_operator}{相关文章}。

\subsection{拉普拉斯算子}
在数学中，拉普拉斯算子（Laplace operator 或 Laplacian）是一种微分算子，由欧几里得空间中标量函数的梯度的散度给出。它通常记作符号$\nabla \cdot \nabla$、$\nabla^{2}$（其中$\nabla$为纳布拉算子，nabla operator），或 $\Delta$。在笛卡尔坐标系中，拉普拉斯算子等于函数关于每个独立变量的二阶偏导数之和；在其他坐标系（如柱坐标系和球坐标系）中，拉普拉斯算子也有相应的有用形式。非正式地说，函数$f$在点$p$的拉普拉斯值$\Delta f(p)$衡量的是：以$p$为中心的小球上$f$的平均值与$f(p)$本身的偏离程度。

拉普拉斯算子以法国数学家皮埃尔–西蒙·拉普拉斯（Pierre-Simon de Laplace，1749–1827）命名。他最早将该算子应用于天体力学研究：给定质量密度分布下的引力势的拉普拉斯值是该密度分布的常数倍。满足拉普拉斯方程$\Delta f = 0$的函数称为调和函数（harmonic functions），它们表示真空区域内可能的引力势。

拉普拉斯算子出现在许多描述物理现象的微分方程中：泊松方程（Poisson’s equation）描述电势与引力势，扩散方程（diffusion equation）描述热流与流体流动，波动方程（wave equation）描述波的传播，而薛定谔方程（Schrödinger equation）则描述量子力学中的波函数。在图像处理与计算机视觉领域中，拉普拉斯算子也被用于多种任务，如斑点检测（blob detection）与边缘检测（edge detection）拉普拉斯算子是最简单的椭圆型算子（elliptic operator），并且在霍奇理论（Hodge theory）及德拉姆上同调（de Rham cohomology）的结果中起着核心作用。
\subsection{定义}
拉普拉斯算子（Laplace operator）是定义在$n$维欧几里得空间中的二阶微分算子，其定义为梯度（$\nabla f$）的散度（$\nabla \cdot$）。因此，若$f$是一个二次可微的实值函数，则$f$的拉普拉斯算子为：
\[
\Delta f = \nabla^{2} f = \nabla \cdot \nabla f~
\]
其中符号 $\nabla$ 正式写作
\[
\nabla = \left( \frac{\partial}{\partial x_{1}}, \ldots, \frac{\partial}{\partial x_{n}} \right)~
\]
显式地说，在笛卡尔坐标系$(x_{1}, \ldots, x_{n})$中，函数$f$的拉普拉斯算子是其关于各坐标的非混合二阶偏导数之和：
\[
\Delta f = \sum_{i=1}^{n} \frac{\partial^{2} f}{\partial x_{i}^{2}}
~
\]
作为一个二阶微分算子，拉普拉斯算子将$C^{k}$类函数映射为$C^{k-2}$类函数（当 $k \ge 2$ 时）。它是一个线性算子 $\Delta : C^{k}(\mathbf{R}^{n}) \to C^{k-2}(\mathbf{R}^{n})$，或更一般地，$\Delta : C^{k}(\Omega) \to C^{k-2}(\Omega)$，其中 $\Omega \subseteq \mathbf{R}^{n}$ 为任意开集。

另一种等价定义为：
\[
\nabla^{2} f(\vec{x})
= \lim_{R \to 0} \frac{2n}{R^{2}} \left(f_{\text{shell}_{R}} - f(\vec{x})\right)
= \lim_{R \to 0} \frac{2n}{A_{n-1} R^{1+n}}
\int_{\text{shell}_{R}} \left( f(\vec{r}) - f(\vec{x}) \right) \, d r^{n-1}~
\]
其中：
\begin{itemize}
  \item $n$ 是空间的维数；
  \item $f_{\text{shell}_{R}}$ 表示 $f$ 在半径为 $R$ 的 $n$ 维球面上的平均值；
  \item $\displaystyle \int_{\text{shell}_{R}} f(\vec{r}) \, d r^{n-1}$ 是 $f$ 在半径为 $R$ 的 $n$ 维球面上的面积积分；
  \item $A_{n-1}$ 是单位 $n$ 维球面的边界超体积（即单位球面的 $(n-1)$ 维面积）\(^\text{[1]}\)。
\end{itemize}
```latex
\subsection{解析型与几何型的拉普拉斯算子}
关于拉普拉斯算子的定义，存在两种相互冲突的约定：
\begin{enumerate}
  \item 解析型（analytic）拉普拉斯算子：在$\mathbb{R}^{n}$中可定义为
  \[
  \Delta = \nabla^{2} = \sum_{j=1}^{n} \left( \frac{\partial}{\partial x_{j}} \right)^{2},~
  \]
  它是负定的（negative-definite），即对于任意非零的光滑紧支函数 $\varphi \in C_{c}^{\infty}(\mathbb{R}^{n})$，都有
  \[
  \int_{\mathbb{R}^{n}} \overline{\varphi(x)} \, \Delta \varphi(x) \, dx
  = - \int_{\mathbb{R}^{n}} |\nabla \varphi(x)|^{2} \, dx < 0~
  \]
  \item 几何型（geometric）拉普拉斯算子,定义为正定形式（positive-definite）：
  \[
  \Delta = -\nabla^{2}
  = - \sum_{j=1}^{n} \left( \frac{\partial}{\partial x_{j}} \right)^{2}.~
  \]
\end{enumerate}
\subsection{动机}
\subsubsection{扩散现象}
在物理学中的扩散理论（diffusion theory）中，拉普拉斯算子自然地出现在平衡态的数学描述中\(^\text{[2]}\)。  具体而言，若$u$表示某种物理量（例如化学浓度）在平衡时的密度，则对于任意光滑区域$V$，若其中不存在源或汇，则该物理量通过边界 $\partial V$（又记作$S$）的净通量为零：
\[
\int_{S} \nabla u \cdot \mathbf{n} \, dS = 0,~
\]
其中$\mathbf{n}$是$V$边界的外向单位法向量。

根据散度定理（divergence theorem），有
\[
\int_{V} \operatorname{div} \nabla u \, dV
= \int_{S} \nabla u \cdot \mathbf{n} \, dS = 0~
\]

由于上述式子对所有光滑区域 $V$ 都成立，因此可推出
\[
\operatorname{div} \nabla u = \Delta u = 0~
\]
该等式左侧的算子$\Delta$即为拉普拉斯算子，而整个方程$\Delta u = 0$被称为拉普拉斯方程（Laplace’s equation）。其解，即拉普拉斯算子恒为零的函数，表示在扩散平衡状态下可能的密度分布。

在非平衡扩散的情况下，拉普拉斯算子本身具有物理意义：它度量了某一点在化学浓度场中作为“源”或“汇”的强度，这一解释由扩散方程给出了严格的数学定义。这种关于拉普拉斯算子的物理解释，也可以从其“平均值性质”进一步理解。
\subsubsection{平均值性质}
设 $f : \mathbb{R}^{n} \to \mathbb{R}$ 为一个二次连续可微函数，$p \in \mathbb{R}^{n}$。以$p$为中心、半径为$h$的球体上$f$的平均值为\(^\text{[3]}\)：
\[
\overline{f}_{B}(p, h)
= f(p) + \frac{\Delta f(p)}{2(n + 2)}h^{2} + o(h^{2}),
\quad \text{当 } h \to 0~
\]
类似地，以$p$为中心、半径为$h$的球面（即球的边界）上$f$的平均值为：
\[
\overline{f}_{S}(p, h)
= f(p) + \frac{\Delta f(p)}{2n}h^{2} + o(h^{2}),
\quad \text{当 } h \to 0~
\]
\subsubsection{与势函数相关的密度}
若 $\varphi$ 表示与电荷分布 $q$ 相对应的静电势（electrostatic potential），则该电荷分布由 $\varphi$ 的拉普拉斯算子的负值给出：
\[
q = -\varepsilon_{0}\,\Delta \varphi,~
\]
其中 $\varepsilon_{0}$ 为真空介电常数（electric constant）。

该结果来源于高斯定律（Gauss’s law）。确实，若 $V$ 是任意光滑区域，其边界为 $\partial V$，则根据高斯定律，静电场 $\mathbf{E}$ 穿过边界的通量与所包围的电荷成正比：
\[
\int_{\partial V} \mathbf{E} \cdot \mathbf{n}\, dS
= \int_{V} \operatorname{div} \mathbf{E}\, dV
= \frac{1}{\varepsilon_{0}} \int_{V} q\, dV,~
\]
其中第一步等号来自散度定理。

由于静电场是势函数的负梯度，即$\mathbf{E} = -\nabla \varphi$，上式可写为：
\[
-\int_{V} \operatorname{div}(\operatorname{grad} \varphi)\, dV
= \frac{1}{\varepsilon_{0}} \int_{V} q\, dV~
\]
由于该等式对所有区域 $V$ 都成立，因此可推出：
\[
\operatorname{div}(\operatorname{grad} \varphi)
= -\frac{1}{\varepsilon_{0}} q~
\]
同理，重力势（gravitational potential）的拉普拉斯算子的负值等于质量分布（mass distribution）。在许多情况下，电荷分布或质量分布$q$已知，而相应的势函数$\varphi$未知。求解满足适当边界条件的势函数等价于求解泊松方程（Poisson’s equation）。
\subsubsection{能量最小化原理}
拉普拉斯算子出现在物理学中的另一个动机是：在区域$U$上满足$\Delta f = 0$的函数，正是使狄利克雷能量泛函（Dirichlet energy functional）处于驻点（stationary）状态的函数：
\[
E(f) = \frac{1}{2} \int_{U} \lVert \nabla f \rVert^{2} \, dx~
\]
为说明这一点，设 $f : U \to \mathbf{R}$ 是一个函数，且 $u : U \to \mathbf{R}$ 是一个在区域边界上取零的函数。则有：
\[
\left. \frac{d}{d\varepsilon} \right|_{\varepsilon = 0} E(f + \varepsilon u)
= \int_{U} \nabla f \cdot \nabla u \, dx
= - \int_{U} u \, \Delta f \, dx~
\]
其中最后一个等号由格林第一恒等式（Green’s first identity）得出。

这一推导表明：若$\Delta f = 0$，则能量泛函$E$在$f$处为驻点。反之，若$E$在$f$附近为驻点，则由变分法基本引理（fundamental lemma of the calculus of variations）可得$\Delta f = 0$。
\subsection{坐标表示}
\subsubsection{二维情形}
在二维情况下，拉普拉斯算子可以表示为：

\textbf{笛卡尔坐标系下：}
\[
\Delta f = \frac{\partial^{2} f}{\partial x^{2}} + \frac{\partial^{2} f}{\partial y^{2}},~
\]
其中$x$与$y$是平面上的标准笛卡尔坐标。

\textbf{极坐标系下}：
\[
\begin{aligned}
\Delta f
&= \frac{1}{r} \frac{\partial}{\partial r} \left( r \frac{\partial f}{\partial r} \right)
    + \frac{1}{r^{2}} \frac{\partial^{2} f}{\partial \theta^{2}} \\[4pt]
&= \frac{\partial^{2} f}{\partial r^{2}}
    + \frac{1}{r} \frac{\partial f}{\partial r}
    + \frac{1}{r^{2}} \frac{\partial^{2} f}{\partial \theta^{2}},
\end{aligned}~
\]
其中$r$表示径向距离，$\theta$表示角度。
\subsubsection{三维情形}
在三维空间中，拉普拉斯算子可根据所采用的坐标系写成不同形式：

\textbf{笛卡尔坐标系}：
\[
\Delta f
= \frac{\partial^{2} f}{\partial x^{2}}
+ \frac{\partial^{2} f}{\partial y^{2}}
+ \frac{\partial^{2} f}{\partial z^{2}}~
\]

\textbf{柱坐标系：}
\[
\Delta f
= \frac{1}{\rho}
  \frac{\partial}{\partial \rho}
  \left( \rho \frac{\partial f}{\partial \rho} \right)
+ \frac{1}{\rho^{2}}
  \frac{\partial^{2} f}{\partial \varphi^{2}}
+ \frac{\partial^{2} f}{\partial z^{2}}~
\]
其中 $\rho$ 表示径向距离，$\varphi$ 表示方位角，$z$ 表示高度。

\textbf{球坐标系：}
\[
\Delta f
= \frac{1}{r^{2}}
  \frac{\partial}{\partial r}
  \left( r^{2} \frac{\partial f}{\partial r} \right)
+ \frac{1}{r^{2} \sin \theta}
  \frac{\partial}{\partial \theta}
  \left( \sin \theta \frac{\partial f}{\partial \theta} \right)
+ \frac{1}{r^{2} \sin^{2} \theta}
  \frac{\partial^{2} f}{\partial \varphi^{2}}~
\]
或等价地写作：
\[
\Delta f
= \frac{1}{r} \frac{\partial^{2}}{\partial r^{2}}(r f)
+ \frac{1}{r^{2} \sin \theta}
  \frac{\partial}{\partial \theta}
  \left( \sin \theta \frac{\partial f}{\partial \theta} \right)
+ \frac{1}{r^{2} \sin^{2} \theta}
  \frac{\partial^{2} f}{\partial \varphi^{2}}~
\]

将前两项展开后，有：
\[
\Delta f
= \frac{\partial^{2} f}{\partial r^{2}}
+ \frac{2}{r} \frac{\partial f}{\partial r}
+ \frac{1}{r^{2} \sin \theta}
  \left( \cos \theta \frac{\partial f}{\partial \theta}
       + \sin \theta \frac{\partial^{2} f}{\partial \theta^{2}} \right)
+ \frac{1}{r^{2} \sin^{2} \theta}
  \frac{\partial^{2} f}{\partial \varphi^{2}}~
\]
其中 $\varphi$ 表示方位角（azimuthal angle），$\theta$ 表示天顶角或余纬角（zenith angle / co-latitude）。

特别地，上式等价于：
\[
\Delta f
= \frac{\partial^{2} f}{\partial r^{2}}
+ \frac{2}{r} \frac{\partial f}{\partial r}
+ \frac{1}{r^{2}} \Delta_{S^{2}} f~
\]
其中 $\Delta_{S^{2}} f$ 为单位球面上的拉普拉斯–贝尔特拉米算子（Laplace–Beltrami operator）。

\textbf{一般曲线坐标系 $(\xi^{1}, \xi^{2}, \xi^{3})$：}
\[
\Delta
= \nabla \xi^{m} \cdot \nabla \xi^{n}
  \frac{\partial^{2}}{\partial \xi^{m} \partial \xi^{n}}
+ \nabla^{2} \xi^{m}
  \frac{\partial}{\partial \xi^{m}}
= g^{mn} \left(
  \frac{\partial^{2}}{\partial \xi^{m} \partial \xi^{n}}
  - \Gamma^{l}_{mn}
    \frac{\partial}{\partial \xi^{l}}
\right)~
\]
其中对重复指标隐含求和；$g^{mn}$为逆度量张量（inverse metric tensor），$\Gamma^{l}_{mn}$为所选坐标系的克里斯托费尔符号（Christoffel symbols）。
\subsubsection{N 维情形}
在$N$维的任意曲线坐标系$(\xi^{1}, \ldots, \xi^{N})$中，拉普拉斯算子可通过逆度量张量$g^{ij}$表示为：
\[
\Delta
= \frac{1}{\sqrt{\det g}}
  \frac{\partial}{\partial \xi^{i}}
  \left(
    \sqrt{\det g}\,
    g^{ij}
    \frac{\partial}{\partial \xi^{j}}
  \right)~
\]
该公式来自于 Voss–Weyl 散度公式\(^\text{[4]}\)。

在$N$维球坐标系中，设参数化为$\mathbf{x} = r \boldsymbol{\theta} \in \mathbf{R}^{N}$，其中$r$表示正实半径，$\boldsymbol{\theta}$ 为单位球面 $S^{N-1}$上的点，则有：
\[
\Delta f
= \frac{\partial^{2} f}{\partial r^{2}}
+ \frac{N - 1}{r} \frac{\partial f}{\partial r}
+ \frac{1}{r^{2}} \Delta_{S^{N-1}} f~
\]
其中 $\Delta_{S^{N-1}}$ 为 $(N - 1)$ 维球面上的拉普拉斯–贝尔特拉米算子（Laplace–Beltrami operator），又称球面拉普拉斯算子（spherical Laplacian）。

上述两个径向导数项也可等价地写作：
\[
\frac{1}{r^{N-1}}
\frac{\partial}{\partial r}
\left(
  r^{N-1} \frac{\partial f}{\partial r}
\right)~
\]
因此，对于定义在$S^{N-1} \subset \mathbf{R}^{N}$ 上的函数，其球面拉普拉斯算子可以通过将该函数扩展到$\mathbf{R}^{N} \setminus \{0\}$上、并保持沿射线方向不变（即为零次齐次函数）的普通拉普拉斯算子来计算。
\subsection{欧几里得不变性}
拉普拉斯算子在所有欧几里得变换下（旋转与平移）保持不变。例如，在二维情形下，对于任意的$\theta, a, b$，都有：
\[
\Delta \bigl( f(x \cos \theta - y \sin \theta + a,\;
x \sin \theta + y \cos \theta + b) \bigr)
= (\Delta f)(x \cos \theta - y \sin \theta + a,\;
x \sin \theta + y \cos \theta + b).~
\]
在任意维度中，若$\rho$表示一个旋转变换，则有：
\[
\Delta (f \circ \rho) = (\Delta f) \circ \rho.~
\]
同理，若 $\tau$ 表示一个平移变换，则：
\[
\Delta (f \circ \tau) = (\Delta f) \circ \tau.~
\]
（更一般地，当$\rho$是任意正交变换——例如反射——时，上述结论仍然成立。）

事实上，所有与欧几里得变换可交换的、具有常系数的标量线性微分算子的代数，
正是由拉普拉斯算子生成的多项式代数。
\subsection{谱理论}
拉普拉斯算子的谱由所有满足下式的特征值$\lambda$构成：
\[
- \Delta f = \lambda f,~
\]
其中$f$是对应的特征函数（eigenfunction）。该方程称为Helmholtz 方程（Helmholtz equation）。

若$\Omega$是$\mathbf{R}^{n}$中的一个有界区域，则拉普拉斯算子的特征函数构成 Hilbert 空间$L^{2}(\Omega)$的一组正交规范基（orthonormal basis）。该结果本质上来源于紧自伴算子谱定理（spectral theorem on compact self-adjoint operators），应用于拉普拉斯算子的逆算子（根据 Poincaré 不等式与 Rellich–Kondrachov 定理可知其为紧算子）\(^\text{[5]}\)。此外，可以证明这些特征函数是无穷可微函数（infinitely differentiable functions）\(^\text{[6]}\)。更一般地，上述结论对任意带边界的紧黎曼流形（Riemannian manifold）上的拉普拉斯–贝尔特拉米算子（Laplace–Beltrami operator）亦成立，或者更广泛地，对定义在有界区域上、具有光滑系数的任意椭圆算子的狄利克雷特特征值问题也成立。当 $\Omega$是$n$维球面$S^{n}$时，拉普拉斯算子的特征函数即为球谐函数（spherical harmonics）。
\subsection{向量拉普拉斯算子}
向量拉普拉斯算子（也记作$\nabla^{2}$是定义在向量场上的一个微分算子 \(^\text{[7]}\)。它与标量拉普拉斯算子相似：标量拉普拉斯作用于标量场并返回标量量，而向量拉普拉斯作用于向量场并返回向量量。在正交笛卡尔坐标系下，向量拉普拉斯算子返回的向量场等于对向量场的每个分量分别施加标量拉普拉斯算子后所得的结果。

设$\mathbf{A}$是一个向量场，则其向量拉普拉斯定义为：
\[
\nabla^{2} \mathbf{A}
= \nabla (\nabla \cdot \mathbf{A})
- \nabla \times (\nabla \times \mathbf{A}).~
\]
此定义可视为Helmholtz 分解（Helmholtz decomposition）在向量拉普拉斯情形下的形式。

\textbf{在笛卡尔坐标系下：}当 $\nabla \cdot \mathbf{A} = 0$ 时，上式化简为更简单的形式：
\[
\nabla^{2} \mathbf{A}
= (\nabla^{2} A_{x},\, \nabla^{2} A_{y},\, \nabla^{2} A_{z})~
\]
其中$A_{x}, A_{y}, A_{z}$为向量场$\mathbf{A}$的各分量，且作用于每个分量左侧的$\nabla^{2}$为标量拉普拉斯算子。这实际上是拉格朗日公式（Lagrange’s formula）的一个特殊情形（参见向量三重积）。

关于其他坐标系中向量拉普拉斯的表达式，参见：柱坐标系与球坐标系中的 $\boldsymbol{\nabla}$ 运算。
\subsubsection{推广}
任意张量场$\mathbf{T}$（“张量”包括标量与向量）的拉普拉斯算子，可定义为该张量梯度的散度：
\[
\nabla^{2} \mathbf{T} = (\nabla \cdot \nabla) \mathbf{T}.~
\]
\textbf{特殊情形：}若$\mathbf{T}$为标量（即 0 阶张量），则拉普拉斯算子退化为熟知的标量形式；若$\mathbf{T}$为向量（即 1 阶张量），其梯度为一个协变导数，得到一个 2 阶张量，而其散度再返回一个向量。上文所给出的向量拉普拉斯公式可用于避免直接使用张量计算，且可以证明它等价于下述关于向量梯度的雅可比矩阵散度形式：
\[
\nabla \mathbf{T}
= (\nabla T_{x}, \nabla T_{y}, \nabla T_{z})
= 
\begin{bmatrix}
T_{xx} & T_{xy} & T_{xz} \\
T_{yx} & T_{yy} & T_{yz} \\
T_{zx} & T_{zy} & T_{zz}
\end{bmatrix},
\quad \text{其中 } T_{uv} \equiv \frac{\partial T_{u}}{\partial v}~
\]
类似地，向量$\mathbf{A}$与另一个向量场$\mathbf{B}$的梯度（一个二阶张量）的点积，可视为矩阵乘积：
\[
\mathbf{A} \cdot \nabla \mathbf{B}
= 
\begin{bmatrix}
A_{x} & A_{y} & A_{z}
\end{bmatrix}
\nabla \mathbf{B}
=
\begin{bmatrix}
\mathbf{A} \cdot \nabla B_{x},
\; \mathbf{A} \cdot \nabla B_{y},
\; \mathbf{A} \cdot \nabla B_{z}
\end{bmatrix}.~
\]
需要注意，该恒等式依赖于坐标系，不具有一般协变性。
